\section*{第 \ref{ch:b2:conjugate-additions} 章——共轭羰基化合物：迈克尔与罗宾逊}

\begin{solution}{exo:b2:conjugate-additions:1}
\begin{itemize}
  \item \ce{CH3MgBr}：主要发生 \omterm{def:b2:conjugate-additions:addition-modes}{1,2-加成}（硬、不可逆），生成 1-甲基环己-2-烯-1-醇。
  \item \ce{(CH3)2CuLi}：1,4-加成，生成 3-甲基环己酮。
  \item 苯硫酚与碱作用：1,4-加成（硫醇盐较软），生成 3-（苯硫基）环己酮。
  \item \ce{LiAlH4}：\omterm{def:b2:conjugate-additions:addition-modes}{1,2-加成}，生成环己-2-烯-1-醇。
\end{itemize}
\end{solution}

\begin{solution}{exo:b2:conjugate-additions:2}
给体：2,4-戊二酮、硝基甲烷、丙二酸二乙酯（具有酸性的 \ce{C-H}）。受体：
丙烯腈、丙烯酸甲酯、丁-3-烯-2-酮（\ce{C=C} 与吸电子基团共轭）。
\end{solution}

\begin{solution}{exo:b2:conjugate-additions:3}
\ce{C4H9Br + 2Li -> C4H9Li + LiBr}，然后 \ce{2C4H9Li + CuI -> (C4H9)2CuLi + LiI}，于乙醚中低
温、惰性气体保护下进行。
\end{solution}

\begin{solution}{exo:b2:conjugate-additions:4}
（3-氧代丁基）丙二酸二乙酯，\ce{(EtOOC)2CHCH2CH2COCH3}。水解得取代丙二酸，
受热脱去 \ce{CO2}（见\cref{prop:b2:enolates-aldol:decarboxylation}）：
5-氧代己酸，\ce{HOOCCH2CH2CH2COCH3}。
\end{solution}

\begin{solution}{exo:b2:conjugate-additions:5}
二甲基铜锂在乙醚中低温反应，随后水溶液后处理：\omterm{def:b2:conjugate-additions:addition-modes}{共轭加成}。
甲基溴化镁较硬，主要发生 \omterm{def:b2:conjugate-additions:addition-modes}{1,2-加成}，
生成主要产物 1-甲基环己-2-烯-1-醇。
\end{solution}

\begin{solution}{exo:b2:conjugate-additions:6}
\ce{Et3N} 夺去少量 \ce{EtSH}，生成 \ce{EtS-}；硫醇盐加成到 $\beta$ 碳原子上，所得
\omterm{def:b2:enolates-aldol:enolate}{烯醇负离子}从另一分子 \ce{EtSH} 获得质子，并再生 \ce{EtS-}。产物：
产物为 4-（乙硫基）丁-2-酮。硫体积大且易极化，孤电子对能量较高，属于\omterm{def:b2:frontier-orbitals:hard-soft}{软
亲核试剂}，受\omterm{def:b2:frontier-orbitals:control}{轨道控制}，故进攻 \omterm{def:b2:frontier-orbitals:frontier}{LUMO} 系数最大处，并且加成是
可逆的，这也有利于更稳定的 1,4-加成产物。
\end{solution}

\begin{solution}{exo:b2:conjugate-additions:7}
氰根更快进攻电正性更强的羰基碳，故氰醇是动力学产物。其
生成是可逆的；温度较高、时间充足时，氰根被释放出来，并在
$\beta$ 碳上重新加成，得到保留较强 \ce{C=O} $\pi$ 键的加成物，即热力学产物。
\end{solution}

\begin{solution}{exo:b2:conjugate-additions:8}
以 2-甲基环己烷-1,3-二酮为原料时得到维兰德–米舍尔酮（\omterm{def:b2:conjugate-additions:robinson}{罗宾逊成环}的示例）。
以环己酮为原料时得到 4,4a,5,6,7,8-六氢萘-2(3\textit{H})-酮，
为双环烯酮，其 \ce{C=C} 终止于一个环稠合碳原子，无角甲基。
\end{solution}

\begin{solution}{exo:b2:conjugate-additions:9}
\ce{CH3COCH2CH2CH2COCH3}：C2 与 C6 为羰基碳，从 C2 到 C6 含两端共五个碳，
是一个 1,5-二酮。断开 C3--C4：给体为丙酮在 C3 上的\omterm{def:b2:enolates-aldol:enolate}{烯醇负离子}（更好的是 3-氧代丁酸乙酯，其
酯基随后经水解与\omterm{def:b2:enolates-aldol:decarboxylation}{脱羧}除去）；受体为丁-3-烯-2-酮。条件：
3-氧代丁酸乙酯、丁-3-烯-2-酮，乙醇中的催化量乙醇钠；随后
用酸性水溶液处理并加热。
\end{solution}

\begin{solution}{exo:b2:conjugate-additions:10}
\omterm{def:b2:enolates-aldol:aldol}{羟醛}反应形成 \ce{C-C} 键，却把 \ce{C=O} $\pi$ 键换成 \ce{C-O} $\sigma$ 键；其
平衡不大，两分子合为一个，故 $K$ 不大，逆\omterm{def:b2:enolates-aldol:aldol}{羟醛}反应容易进行。
\omterm{def:b2:conjugate-additions:michael}{迈克尔加成}形成 \ce{C-C} 键，代价是较弱的 \ce{C=C} $\pi$ 键，并
保留两个 \ce{C=O}：该过程明显有利。逆反应需由酮的\omterm{def:b2:enolates-aldol:enolate}{烯醇负离子}排出稳定的丙二酸根负离子，
平衡强烈偏向加成物一侧，以致催化量碱存在下加成物得以保留。
\end{solution}

\begin{solution}{exo:b2:conjugate-additions:11}
\ce{C(COOEt)2(CH2CH2COOCH3)2}。第一次\omterm{def:b2:conjugate-additions:michael}{迈克尔加成}后，中心碳上仍有一个
氢位于两个酯基之间，酸性不变：碱夺去它，随即发生第二次加成。
\end{solution}

\begin{solution}{exo:b2:conjugate-additions:12}
经取代程度较高的\omterm{def:b2:enolates-aldol:enolate}{烯醇负离子}（朝向带甲基的碳）反应，得到
甲基位于环稠合碳上的八氢萘酮，即 4a-甲基-4,4a,5,6,7,8-六氢萘-2(3\textit{H})-酮。经
取代程度较低的\omterm{def:b2:enolates-aldol:enolate}{烯醇负离子}（在 \ce{CH2} 一侧）进行时，甲基位于紧邻
稠合碳的环碳上。前者在平衡条件下更有利（醇中的醇盐，或\omterm{def:b2:amines:class}{仲胺}与酸的\omterm{def:b2:amines:imine}{烯胺}
途径），这些条件给出热力学控制、取代程度较高的\omterm{def:b2:enolates-aldol:enolate}{烯醇负离子}
（见\cref{prop:b2:enolates-aldol:enolate-control}）。
\end{solution}

\begin{solution}{pb:b2:conjugate-additions:1}
\textbf{1.} 环己烷环，C1 与 C3 为 \ce{C=O}，C2 上连有甲基；C2 上的氢
酸性最强。
\textbf{2.} 失去该氢得到负离子，其电荷离域在 C2 与两个氧原子上；对
环己酮只有一个氧参与分担电荷，其酸性弱到无法在水中测定。
\textbf{3.} 电荷位于 C2 上；\ce{C1=C2}，同时 C1 上带 \ce{O^-}；\ce{C2=C3}，同时 C3 上带 \ce{O^-}。
\textbf{4.} 能：以氢氧根计，$K = 10^{14.0 - 5.3} = 10^{8.7}$（水作为共轭酸，由
$K_e$ 求得）；用乙醇钠时，$10^{15.9-5.3} = 10^{10.6}$。两种情况均进行完全。
% ledger: ke:25C, pka:ethanol
\textbf{5.} \ce{C1(OH)=C2(CH3)}：C2 上仍有一个氢，\omterm{def:b2:enolates-aldol:tautomerism}{烯醇}只需要这一个氢；甲基
仅取代了第二个氢。
\textbf{6.} 软：电荷分散在三个原子上，负离子得到稳定。它进攻
$\beta$ 碳（见\cref{prop:b2:conjugate-additions:regio}）。
\textbf{7.} \omterm{def:b2:enolates-aldol:enolate}{烯醇负离子}的 C2 与丁-3-烯-2-酮的 \ce{CH2} 端成键；生成的\omterm{def:b2:enolates-aldol:enolate}{烯醇负离子}
获得一个质子，得到 2-甲基-2-（3-氧代丁基）环己烷-1,3-二酮。
\textbf{8.} 每个新生成的\omterm{def:b2:enolates-aldol:enolate}{烯醇负离子}从一分子二酮（或溶剂）获得质子，
从而再生给体\omterm{def:b2:enolates-aldol:enolate}{烯醇负离子}；碱并未消耗。
\textbf{9.} \omterm{def:b2:enolates-aldol:enolate}{烯醇负离子}是两可的；与软的碳亲电试剂作用、受\omterm{def:b2:frontier-orbitals:control}{轨道控制}时，它
在碳上反应，\ce{C-C} 加成物保留两个 \ce{C=O}，这一组合更稳定；
\ce{O}-加成物若形成，则会逆转。
\textbf{10.} 一个：C2，与 C1、C3、甲基和侧链相连。
\textbf{11.} 侧链羰基碳与环上羰基 C1（或 C3）：\ce{C(=O)}--\ce{CH2}--\ce{CH2}--C2--C1，
含两端共五个碳。
\textbf{12.} 为使较贵的二酮完全反应，丁-3-烯-2-酮只能略微过量，
因其毒性很大且易聚合。
\textbf{13.} 环上 C4 与 C6（分别紧邻 C3 与 C1）、紧邻羰基的侧链 \ce{CH2}，以及侧链
末端 \ce{CH3}。C2 上没有。
\textbf{14.} 末端 \ce{CH3} 的\omterm{def:b2:enolates-aldol:enolate}{烯醇负离子}进攻环上羰基，所成环包含该
碳、侧链羰基碳、两个 \ce{CH2}、C2 与被进攻的羰基碳，共六个原子。
\textbf{15.} 侧链内部 \ce{CH2} 的\omterm{def:b2:enolates-aldol:enolate}{烯醇负离子}将形成四元环，张力较大。环上
\omterm{def:b2:enolates-aldol:enolate}{烯醇负离子}（C4 或 C6）进攻侧链羰基，同样形成六元环，但为桥环，该\omterm{def:b2:enolates-aldol:aldol}{羟醛}
不能脱水（双键将位于小双环体系的桥头碳上），并逆向转化，
而稠合\omterm{def:b2:enolates-aldol:aldol}{羟醛}借脱水得以转化。
\textbf{16.} \omterm{def:b2:enolates-aldol:aldol}{羟醛}在原环羰基碳上带有 \ce{OH}，该碳现位于环稠合处；失水
时与原 \ce{CH3} 碳上一个氢同时消去（该碳现位于侧链酮的 $\alpha$ 位），得到
与该酮共轭的 \ce{C=C}。
\textbf{17.} 新生成的双键与羰基共轭：经\omterm{def:b2:enolates-aldol:enolate}{烯醇负离子}进行的 E1cB 消除，
加热有利于其进行（见\cref{prop:b2:enolates-aldol:aldol-equilibrium}）。
\textbf{18.} 一种饱和酮和一种 $\alpha,\beta$-不饱和酮（环己烯酮）。
\textbf{19.} 带甲基的环稠合碳：连有四个不同的取代基。
\textbf{20.} 否：二酮是非手性的，其两个羰基被进攻的机会相同（二者相对于侧链互为镜像
关系）；与非手性试剂反应，产物为外消旋体。手性胺
催化剂可偏向其中一个羰基，给出过量的一种对映异构体（见本科三年级分册）。
\textbf{21.} 二酮 \ce{C7H10O2}：\qty{126.155}{g/mol}；\ce{C4H6O}：\qty{70.091}{g/mol}；\ce{C11H14O2}：
\qty{178.231}{g/mol}。
\textbf{22.} \ce{C7H10O2 + C4H6O -> C11H14O2 + H2O}：两侧均为 C 11、H 16、O 3。
\textbf{23.} $10.0/126.155 = \qty{0.07927}{mol} = \qty{79.3}{mmol}$。
\textbf{24.} 所需维兰德–米舍尔酮的质量约为
$0.07927 \times 178.231 \times 0.70 = \boldsymbol{\approx \qty{9.89}{g}}$。
\end{solution}
